\chapter*{Ejercicios matemáticos de la guía}
\addcontentsline{toc}{chapter}{Ejercicios matemáticos de la guía}
\noindent Estos ejercicios consolidan resultados de la clase y corresponden al
trabajo autónomo, no al tiempo presencial. Las soluciones justificadas están
separadas en \texttt{solej01.tex}.
\begin{enumerate}
\item \textbf{[Derivación]} Demuestre la regla de la cadena para $T=4$.
Para $X_1,X_2$ Bernoulli$(1/2)$ independientes y $X_3=X_1$,
contraste $p(1,0,1)$ con la aproximación de orden uno.
\emph{Autoverificación:} $1/4$ frente a $1/8$.
\item \textbf{[Cálculo]} Para $c=(6,3,1,0,0)$ y $K=5$,
calcule la predictiva usando dos \emph{priors} simétricas:
Dirichlet con parámetros unitarios y Dirichlet con parámetros $1/2$.
Compare la masa sobre categorías no vistas.
\emph{Autoverificación:} $1/15$ frente a $1/25$ en cada categoría nula.
\item \textbf{[Demostración]} Pruebe Gibbs desde $\ln t\le t-1$,
justifique cuándo se da la igualdad y deduzca
$H(p,q)=H(p)+D_{\rm KL}(p\|q)$.
\emph{Autoverificación:} KL es no negativa, pero generalmente asimétrica.
\item \textbf{[Cálculo]} Un predictor asigna a cuatro resultados observados
$(1/2,1/4,1/8,1/2)$. Calcule pérdida total, pérdida media y
PPL. Repita cuando la tercera probabilidad sea $1/16$.
\emph{Autoverificación:} PPL inicial $2^{7/4}$, final $4$.
\item \textbf{[Aplicación]} En un mismo texto de 100 caracteres, A
produce 100 tokens con PPL 3.2 y B 30 tokens con PPL 25.
Convierta a bits totales y BPC y enumere cuatro condiciones necesarias
para interpretar una diferencia.
\emph{Autoverificación:} BPC(A) $\approx1.6781$;
BPC(B) $\approx1.3932$.
\end{enumerate}
\section*{Problemas suplementarios}
\begin{enumerate}
\item[(S1.A)] Para una cadena de Markov homogénea, relacione su
entropía condicional $H(X_t\mid X_{t-1})$ con las tasas de
codificación cuando existe una distribución estacionaria.
\item[(S1.B)] Formalice por qué una frecuencia bruta de pares solapados
no es la reducción de longitud de una fusión BPE.
\item[(S1.C)] \emph{Semilla IC10:} formule una hipótesis falsable que
compare tasa de segmentación en corpus general y técnico, e indique
qué control haría falta para atribuir un efecto al tokenizador.
\end{enumerate}
